% !TeX root = ../papers/nc-p3.tex
\section{有限辺の多重点条件と二次元単純から作る六次行列式}
\label{sec:finite-edge-symbolic-powers}
\begin{theorem}\label{thm:finite-edge-symbolic-powers}
有限辺の長さ二十の点スキーム $D\subset\mathbb P^3$ に対し、全ての $m>0$ で
$H^0(\mathbb P^3,\mathcal I_D^m(2m))=0$ である。ここで冪はイデアル層の積で、
$D$ が被約なら各点に重複度 $m$ を課す通常の多重点条件である。
また長さ4で $D$ を切る直線は互いに交わらず、高々二十本である。
\end{theorem}
\begin{proof}
一辺の六関係の空間を $R$、その同時評価を
$f:\mathbb P^3\times\mathbb P^3\dashrightarrow\mathbb P(R^*)$ とする。
基底スキームは指定された長さ二十のグラフである。左点 $p$ を固定したときの全評価の
線型包を $W_p\subset R^*$ とする。$p\notin D$ なら四次元、$p\in D$ の閉点なら三次元である。
$D$ の局所イデアルはこの六行四列評価行列の最大小行列で生成される。
階数3の点で単元小行列により右点を一意に解けば残り三式がちょうど局所余因子だからである。

直線 $L$ 上で $W_p$ が一般に一定なら、四次の最大小行列は全て同じ非零四次式の倍であり、
$L\cap D$ は長さ4である。逆も制限四次式の次数から従う。
一定の $W$ の零化二平面 $K_L\subset R$ の関係は $L\times\mathbb P^3$ 全体で消える。
そのような二直線が交われば交点で二つの $K_L$ は高々三次元の評価核に入るため
共通の非零関係を持つ。この関係は両直線の張る平面を消し、階数1となる。
関係に純テンソルがないことに反する。従って直線は互いに素で、各々が有限な台に会うから有限である。

任意の整曲面 $V\subset\mathbb P^3$ で $f(V\times\mathbb P^3)$ が稠密になることを示す。
稠密でないとし、その閉包を $Y$ とする。$p\in V$ 一般で三平面 $\mathbb P(W_p)$ を含む。
像次元が3なら全て同じ三平面となる。するとその零化関係は $V$ を消し、非零関係の左核の
射影次元が高々1であることに反する。よって $Y$ は四次元超曲面である。
一般の滑らかな $y\in Y$ における接超平面の方程式を $A_y\in R$ とする。
$y\in\mathbb P(W_p)$ ならこの三平面全体が接超平面に含まれ、$p$ は $A_y$ の左核に入る。
一般繊維は曲線であり、$p\notin D$ では右点は $p,y$ から一意に解ける。
従って一般繊維の各一般点の近くで $p$ はこの左核の直線 $L_y$ を動く。
$A_y$ の階数はちょうど2で、$L_y\subset V$ である。

$V$ が平面でなければ、一般の滑らかな $p\in V$ を通る $V$ 上の直線は有限本である。
無限なら接平面内のそれらの和が平面を埋めるからである。
一般の $p$ を固定し $y$ を $\mathbb P(W_p)$ で動かすと、一般の $y$ に対する $L_y$ は
このうちの一本に固定される。その各一般 $p'\in L_y$ で一般の $y$ が $\mathbb P(W_{p'})$ に入る。
実際各 $y$ の繊維の左射影はこの直線に稠密で、所属は $D$ の外で閉条件である。
よって直線上で $W_{p'}$ は一定である。先に得た有限本の直線で $V$ を覆えず矛盾する。

$V=\Pi$ が平面であるとする。六関係の制限はなお独立である。零に制限される関係は階数1だからである。
$A_y$ の $\Pi$ への制限は $L_y\subset\Pi$ と階数2から非零の階数1であり、$u\otimes v$ と書く。
一般の対 $(p,q)\in\Pi\times\mathbb P^3$ ではこの接超平面の引き戻しの一次部が消える。
従って $u(p)=v(q)=0$ である。これは一方の因子だけ動かして微分しても分かる。
一つの $u\otimes v$ に対するこの台は次元3なので、制限六空間内の可分解テンソルには
次元2以上の整な族 $F$ が必要となる。

両因子の線型包のベクトル次元を $l,n$、テンソルの線型包の次元を $b$ とする。
$F$ 上で二つの座標系を掛けた切断空間の次元が $b$ である。
非零切断同士の積は零にならず、射影Segreとの次元比較で $b\ge l+n-1$。
$l<3$ なら左 $l$ 座標を消した部分と右全空間の積（次元 $5-l$）では $b$ 本の式は消える。
残る $6-b$ 本を課しても基底が有限であるため $b\le l+1$、従って $n\le2$。
これは $\dim F\ge2$ により $l=n=2$ で $F$ が二直線の全積となる場合しかないが、
その包は $b=4>l+1$ である。$n<4$ も右側の同じ議論で $l\le2$ となり不可能である。

よって $l=3,n=4$。この場合 $\dim F\ge2$ なら $b\ge7$ である。
両像が曲線なら $F$ は二像の全積で $b=12$。そうでなければ次元2以上の像を持つ方から一般超平面を引き戻す。
その因子はBertiniにより整で、この側の切断の制限核は指定切断の一倍だけである
（整な射影多様体の大域関数は定数）。他側の像はこの因子上で定数でない。
実際一般の二点を通る超平面は存在し、一般超平面上の二点の接続は一般の二点対を支配するので、
他側の二像も一般に異なる。超平面上の積空間の次元は先のSegre評価で、切った側の
元の切断数以上である。制限核には切った式と他側全切断の積が入るため、合わせて $l+n=7$ を得る。
六次元内の族に反し、平面の場合も支配性が示された。

$2m$ 次の非零式が $\mathcal I_D^m$ にあると仮定する。六関係の一般の五線型結合を取り、
二射影空間の積に曲線を切る。基底以外ではBertiniを使え、基底は零次元の六式完全交叉なので
五式は正則列である。残る一式はこの曲線上でグラフのCartier因子を切り、その長さは20である。
左射影の超平面次数は $(h_1+h_2)^5 h_1=10$（点類の係数）。
指定式はこの曲線で $m$ 倍のグラフ因子上に消え、零点次数 $20m$ を全て使う。
この交叉が固有であるのは、各曲面成分と右空間の積が先の支配性により一般の $y$ 上で
有限繊維を持ち、基底は有限であるからである。
一方で同じ支配性によりその一般 $y$ の繊維に基底以外の零点も必ず存在する。矛盾である。
この次数比較はグラフの非被約性も許す。
\end{proof}

\begin{corollary}\label{cor:binary-simple-sextic-determinant}
任意の成分の有限四辺の標準正則窓で、四隅が四行列の恒等式を満たすとする。
二次元均衡単純 $S$ の作用を $A_i$ とする。一つの巡回位置 $j$ で、文字 $r\in V_{j+3}^*$ について
$w(-,A_{j+1},-,r)$ を六行六列に圧縮する。具体的には入力 $V_j^*\otimes S_{j+1}$ を
$A_j(S_j)$ で割り、出力を $V_{j+2}\otimes S_{j+2}\to S_{j+3}$ の核に取る。
その六次行列式がある $S,j$ で零多項式でなければ、窓は左右Noetherかつ中心有限な合成整域である。
逆に中心有限でない正則窓では、この行列式は全単純 $2\delta$ と全四位置で恒等的に零となる。
\end{corollary}
\begin{proof}
二つの二次関係により圧縮写像は定義される。補題\ref{lem:periodic-tail-vertex-bridges}の向きで
$a_i>0$ がどこかにあり、逆向きは消えているとする。各 $D$ の局所Artin環 $C$ で
普遍点を用いて同補題の薄い周期尾部 $P_j$ を取り直す。四移動はスキーム上で恒等であり、
各直線を自明化すれば四周期にできる。$B\otimes C$ 上の点の射影分解は長さ3で項数は $1,3,3,1$。
微分は点の最初・最後の消滅超平面と中央の $w(-,-,-,r)$ である。実際閉点では
Koszul頂点分解とのTor計算による線型分解で、初めの二写像と三次重なりは
定理\ref{thm:finite-cut-point-pairings}の三段と一致する。$C$ 上でも同じ縮約で複体となり、
次数別に自由な項と $C$ 上自由な点加群について閉点の完全性から完全である。

$S\otimes C$ へのHomでは両端の微分はそれぞれ単射と全射で、$C$ 上で分裂する。
閉点で三文字の超平面がある非零ベクトルを全て消せばそのベクトルの全四文字による像が一次元以下となり、
有限辺の非薄単純の評価条件に反する。逆端も双対で同じである。
従って $E_j^3=0$ で、$E_j^2$ は中央写像を両端で簡約した四行四列の余核である。
$E_j^1\to E_{j+1}^1$ は補題と同じ周期完全列で単射、有限長により同型になる。
この完全列の次の部分は
\[
  0\to C^{a_j}\to E_j^2\to E_{j+1}^2\to C^{a_j}\to0 .
\]
端項は体上の $\operatorname{Ext}^{2,3}(k_j,S)$ を $C$ に拡大したものである。
局所完全交叉点の $C$ はArtin Gorensteinで自己入射的、また自由加群は射影的だから
両端は分裂する。よって全ての $E_j^2$ は $C$ 加群として同型であり、少なくとも一つの
自由直和因子を持つ。閉点での階数低下だけではこの結論の代りにならない。

六行六列の圧縮を $C$ に制限した余核には、この $E_j^2$ に加えて自由階数2の商がある。
実際 $r$ の前点によって $w$ の第 $j+2$ 因子はその消滅超平面 $H_{j+2}$ に入る。
圧縮出力から $(V_{j+2}/H_{j+2})\otimes S_{j+2}$ への評価は全射である。
その核が $H_{j+2}\otimes S_{j+2}\to S_{j+3}$ の核であるのは、上で確認した終端の全射性による。
入力側では $w$ の第 $j$ 因子が $H_j$ に入るので $V_j^*\to H_j^*$ を通り、
$A_j(S_j)$ の像を除けば先の四行四列の入力となる。
従って六行列の余核は $C^3$ に全射する。この三つの独立な行を周囲の局所環に持ち上げれば、
可逆な行変換後の三行が全て $\mathcal I_D$ に入る。ゆえに行列式は $\mathcal I_D^3(6)$ に入り、
前定理で零である。反対の向きでは双対単純と逆の窓で同じ論証を使う。
圧縮は逆順で転置となり、同じ行列式の消滅を与える。
定理\ref{thm:all-three-matrix-clifford}により直交する二次元単純一つで全窓が中心有限になるので結論を得る。
\end{proof}
